% DRAFT (2026-09-24): written by Claude from the slides and the lecture notes;
% not yet reviewed by the author.
\section{Intelligent agents}\label{sec:agents}

% ── slides L2 (02_Intelligent_Agents.pdf); AIMA chapter 2 ──
The course builds AI systems as agents: systems that perceive an environment and
act on it. This section fixes the vocabulary for the rest of the course. We
define what an agent is, when its behavior counts as rational, how to describe
the task it faces and its environment, which designs of agent exist, and how an
agent can represent what it knows.

\subsection{Agents and environments}\label{sec:agents-env}

% ── slides L2 p.2–5 ──
To reason about any AI system we first separate the part that decides from the
world it decides about.

\begin{definition}[Agent]\label{def:agent}
  An \textbf{agent} is anything that can act. It perceives its
  \textbf{environment} through \textbf{sensors} and acts on it through
  \textbf{actuators}.
\end{definition}

Where the agent ends and the environment begins is a choice of the designer,
and the agent must be delimited somehow before we can study it. What the sensors
deliver at one instant is a \textbf{percept}. Time runs in discrete steps
$t = 1, 2, \dots$, and at each step the agent receives one percept, so its
experience accumulates as a sequence of percepts.

\begin{definition}[Agent function]\label{def:agent-function}
  The \textbf{agent function}, or \textbf{policy}, is a map
  \[
    f \colon P^* \to A,
  \]
  where $P^*$ is the set of percept histories, the finite sequences of
  percepts, and $A$ is the set of the agent's actions.
\end{definition}

In words: the agent function says which action the agent takes after every
possible history of percepts. We can implement $f$ physically, in a robot, or
simulate it, in a program.

\begin{example}[Vacuum-cleaner world]\label{ex:vacuum-world}
  This worked example runs through the rest of the section. The world has two
  squares, \textsf{A} and \textsf{B}, and the agent is a vacuum cleaner.
  % The slide writes the squares A and B, clashing with the action set A.
  % Squares are set in sans serif here to keep A for the actions.
  \begin{itemize}
    \item The actions are $A = \{\mathit{Left}, \mathit{Right},
      \mathit{Suck}, \mathit{NoOp}\}$.
    \item The position of the agent is $x \in \{\textsf{A}, \textsf{B}\}$.
    \item The status of the square that the agent occupies is
      $s_t \in S = \{\mathit{clean}, \mathit{dirty}\}$. The agent knows it
      only through perception.
    \item The percept at time $t$ is the pair $(x_t, s_t)$.
    \item The task is to clean the environment.
  \end{itemize}
\end{example}

Building an agent means building $f$. \Cref{alg:reflex-vacuum} is a first
attempt: it chooses from the current percept alone. It learns nothing and
contains no AI, because the designer wrote down every case and its action by
hand.

\begin{algorithm}[htbp]
  \caption{Reflex vacuum agent}\label{alg:reflex-vacuum}
  \begin{algorithmic}[1]
    \Require the percept $(\mathit{location}, \mathit{status})$
    \Ensure an action
    \Function{Reflex-Vacuum-Agent}{$\mathit{location}, \mathit{status}$}
      \If{$\mathit{status} = \mathit{dirty}$}
        \State \Return $\mathit{Suck}$
      \ElsIf{$\mathit{location} = \textsf{A}$}
        \State \Return $\mathit{Right}$
      \ElsIf{$\mathit{location} = \textsf{B}$}
        \State \Return $\mathit{Left}$
      \EndIf
    \EndFunction
  \end{algorithmic}
\end{algorithm}

\subsection{Rationality}\label{sec:agents-rationality}

% ── slides L2 p.6–8 ──
Intuitively, \cref{alg:reflex-vacuum} makes sense. To say that it is a
\emph{good} agent, though, we need a criterion. Most definitions of AI ask the
agent to do ``the right thing''. They do not mention intelligence at all:
this narrows the scope of the field, but it makes the goal precise enough to
build agents toward it.

Moral philosophy offers several notions of the right thing. AI adopts
\textbf{consequentialism}: we evaluate an agent by the consequences of its
actions. What counts is the sequence of environment states that its behavior
produces, not the agent's intentions or its own opinion of its success. The
choice is natural for a machine, because consequences can be observed and
measured while intentions cannot. It also comes with a lot of work, since we
must now say which consequences are good. In a simple world we could list
them. As the environment grows more complex, no complete evaluation of every
behavior and its results is possible, so we summarize the good consequences in
a single number.

\begin{definition}[Performance measure]\label{def:performance-measure}
  A \textbf{performance measure} evaluates the goodness of an environment state,
  or of a sequence of environment states.
\end{definition}

\begin{definition}[Rational agent]\label{def:rational-agent}
  For each possible percept sequence, a \textbf{rational agent} selects an
  action that is expected to maximize its performance measure, given the
  evidence provided by the percept sequence and whatever built-in knowledge the
  agent has.
\end{definition}

In words: a rational agent does the best it can with what it has perceived and
what it knew from the start, where ``best'' is fixed by the performance
measure. The summary is where the difficulty lies, as Goodhart's law warns:
``when a measure becomes a target, it ceases to be a good measure''.

\begin{example}[A measure that backfires]\label{ex:goodhart-vacuum}
  This example shows a measure failing. Suppose that we cannot enter the room
  to check whether it is clean, so we reward the vacuum cleaner with the grams
  of dirt that it picks up over a one-hour shift. An agent that maximizes this
  measure can dump the dirt back on the floor and pick it up again: its score
  keeps growing and the room is never cleaner. The measure rewarded a behavior
  that we expected to be useful, instead of the state of the environment that
  we actually wanted.
\end{example}

Rationality also asks less than perfection. A rational agent maximizes
\emph{expected} performance given the percept history, so:
\begin{itemize}
  \item \textbf{rational is not omniscient}: the percepts may not supply all
    the relevant information;
  \item \textbf{rational is not clairvoyant}: the agent has no access to future
    percepts;
  \item \textbf{rational is not successful}: the outcome of an action may not be
    the expected one.
\end{itemize}
In every case, careful engineering of the performance measure, and control of
the behavior that the agent learns, are key.

\subsection{Fixed action patterns}\label{sec:agents-fixed-patterns}

% ── slides L2 p.9 ──
An environment that is completely observable and predictable never changes its
optimal solution, and so it encourages the agent to learn one fixed behavior.
The result is a fragile agent that follows a \textbf{fixed action pattern}: it
keeps executing its policy even when the world changes under it.

\begin{example}[The Sphex wasp]\label{ex:sphex}
  Nature provides a contrast case. After digging a burrow, the Sphex wasp goes
  out, stings a caterpillar and drags it to the burrow. It enters the burrow to
  check that all is well, drags the caterpillar inside and lays its eggs. If an
  entomologist moves the caterpillar a few inches away while the wasp is
  checking the burrow, the wasp goes back to the ``drag the caterpillar'' step
  and continues its plan without modification, checking the burrow again. It
  does so even after dozens of such interventions.
\end{example}

\subsection{Task environments: PEAS}\label{sec:agents-peas}

% ── slides L2 p.10–11 ──
Before designing an agent we describe the task that it must solve. A
\textbf{PEAS} description lists four elements: the \textbf{P}erformance
measure, the \textbf{E}nvironment, the \textbf{A}ctuators and the
\textbf{S}ensors. The actuators and the sensors form the \textbf{agent
architecture}, and the agent is the architecture together with the agent
function.

\begin{example}[Autonomous taxi]\label{ex:peas-taxi}
  \Cref{tab:peas-taxi} is a worked PEAS description.
\end{example}

\begin{table}[htbp]
  \centering
  \renewcommand{\arraystretch}{1.08}
  \caption{PEAS description of an autonomous taxi service.}\label{tab:peas-taxi}
  \begin{tabularx}{\linewidth}{@{}lX@{}}
    \toprule
    \textbf{Element} & \textbf{Autonomous taxi} \\
    \midrule
    Performance measure & safety, reaching the destination, profits, legality,
      comfort, \dots \\
    Environment & streets and freeways, traffic, pedestrians, weather, \dots \\
    Actuators & steering, accelerator, brake, horn, speaker or display, \dots \\
    Sensors & cameras, accelerometers, gauges, engine sensors, keyboard,
      GPS, \dots \\
    \bottomrule
  \end{tabularx}
\end{table}

The more restricted the environment, the easier the design of the agent.

\subsection{Properties of environments}\label{sec:agents-properties}

% ── slides L2 p.12–18. Only the dimensions on the slides: the lecturer said
% that the other dimensions in AIMA (episodic, static, ...) are not required,
% although the summary slide p.32 still names episodic and static. ──
Which agent to build depends on the environment, so we define the environment
first. We classify it along four dimensions, plus a fifth that concerns what
the agent knows.
\begin{itemize}
  \item \textbf{Observable.} The environment is \textbf{fully observable} if the
    sensors give access to its complete state, \textbf{partially observable} if
    faulty sensors or missing information hide part of it, and
    \textbf{unobservable} if the agent has no sensor information at all.
  \item \textbf{Deterministic.} The environment is \textbf{deterministic} if the
    current state and the action uniquely determine the next state, and
    \textbf{nondeterministic} (or stochastic) if the same state and action can
    lead to different next states. Partial observability is one possible cause
    of nondeterminism.
  \item \textbf{Discrete.} The state of the environment, time, the percepts and
    the actions can each take \textbf{discrete} or \textbf{continuous} values.
  \item \textbf{Single agent.} The environment is \textbf{multi-agent} if the
    performance of our agent depends on the actions of other agents, and
    \textbf{single-agent} otherwise.
  \item \textbf{Known.} In a \textbf{known} environment the agent knows its rules,
    that is, which next state follows from each state and action. A known
    environment can still be partially observable. In an \textbf{unknown}
    environment the agent may try to learn the rules.
\end{itemize}

\Cref{tab:env-examples} classifies three environments. The real world sits in
the hardest column: like the taxi, it is partially observable,
nondeterministic, continuous and multi-agent.

\begin{table}[htbp]
  \centering
  \renewcommand{\arraystretch}{1.08}
  \caption{Three environments along the four dimensions.}\label{tab:env-examples}
  \begin{tabular}{@{}lccc@{}}
    \toprule
    \textbf{Characteristic} & \textbf{Sudoku} & \textbf{Backgammon} &
      \textbf{Taxi service} \\
    \midrule
    Fully observable & yes & yes & no \\
    Deterministic    & yes & no  & no \\
    Discrete         & yes & yes & no \\
    Single-agent     & yes & no  & no \\
    \bottomrule
  \end{tabular}
\end{table}

\subsection{Types of agent}\label{sec:agents-types}

% ── slides L2 p.19–28 ──
Agents differ in how they compute the agent function. We go from the simplest
design to the most flexible one; each design adds a component that removes a
limit of the previous one.

\subsubsection{Table-driven agents}\label{sec:agents-table}

The most direct way to build $f$ is to write it down. A
\textbf{table-driven agent} stores its whole percept history and looks up the
action in a table indexed by percept sequences (\cref{alg:table-driven}).

\begin{algorithm}[htbp]
  \caption{Table-driven agent}\label{alg:table-driven}
  \begin{algorithmic}[1]
    \Require a percept
    \Ensure an action
    \Statex \textbf{persistent:} $\mathit{percepts}$, a sequence, initially
      empty; $\mathit{table}$, a table of actions indexed by percept sequences,
      initially fully specified
    \Function{Table-Driven-Agent}{$\mathit{percept}$}
      \State append $\mathit{percept}$ to the end of $\mathit{percepts}$
      \State $\mathit{action} \gets \Call{Lookup}{\mathit{percepts},
        \mathit{table}}$
      \State \Return $\mathit{action}$
    \EndFunction
  \end{algorithmic}
\end{algorithm}

The design is only feasible in simple environments. With a set $P$ of possible
percepts, the table needs one entry for each percept sequence, that is,
$\sum_{t=1}^{T} \abs{P}^t$ entries for an agent that lives $T$ steps. The
complexity of the environment and of the agent makes the table far too large:
the design trades efficiency for effectiveness.

\subsubsection{Simple reflex agents}\label{sec:agents-reflex}

A \textbf{simple reflex agent} chooses its action from the current percept
alone, through \textbf{condition--action rules}: if the percept satisfies a
condition, then the agent takes the corresponding action. The name comes from
the body: like a reflex, the agent reacts to the stimulus without consulting
any memory. \Cref{alg:reflex-vacuum} is a simple reflex agent.

The simplification cuts the cost dramatically. The table of a table-driven agent
grows exponentially with the lifetime $T$, while a simple reflex agent needs one
rule per percept, and its work over $T$ steps grows linearly with $T$. It
simplifies an already naive approach, though: it only works in simple
environments, and small, fully observable environments are rare.

\subsubsection{Model-based agents}\label{sec:agents-model}

A simple reflex agent fails as soon as the current percept does not reveal the
state. A \textbf{model-based agent} keeps an \textbf{internal state}, which
depends on the percept history, and uses it to estimate the current state of
the environment. Two models update the internal state:
\begin{itemize}
  \item the \textbf{transition model} describes how the world evolves, and in
    particular the consequences of the agent's actions;
  \item the \textbf{sensor model} describes how the state of the world shows up
    in the agent's percepts.
\end{itemize}
By approximating the entire state, a model-based agent handles partially
observable environments well. Acquiring such a model is still hard: the lecture
notes that large language models struggle to acquire a model of the world.%
\footnote{The slide points to M.~Mitchell, ``LLMs and world models, part 1''.}

\subsubsection{Goal-based agents}\label{sec:agents-goal}

Knowing the state is not enough to know what to do. A \textbf{goal-based agent}
also represents \textbf{goals}, descriptions of the states it wants to reach,
and it predicts what the world will be like if it takes an action. To choose,
it must weigh the immediate reward of an action against the expected long-term
reward. Two families of methods explore the options: \textbf{search}, which
navigates the tree of possible decisions, and \textbf{planning}, which exploits
knowledge about the world to execute the best actions. Later blocks of the
course study both.

The comparison with a reflex agent shows what goals buy. A reflex agent has no
notion of a goal: its condition--action rules model a single problem, so a new
goal means rebuilding the agent from scratch. A goal-based agent represents
its goals explicitly and can change them.

\subsubsection{Utility-based agents}\label{sec:agents-utility}

A goal only separates success from failure. A \textbf{utility-based agent}
assigns a value, its \textbf{utility}, to every intermediate state or sequence
of states, not only to the final goal: utility works as an internal performance
measure. Utility-based agents are often used in stochastic, partially
observable environments. There the agent maximizes, for each action, the
\textbf{expected utility} over the possible results $i$ of the action:
\[
  \mathrm{EU} = \sum_{i} p_i \, u_i,
\]
where $p_i$ is the probability of result $i$ and $u_i$ is its utility.

\subsubsection{Learning agents}\label{sec:agents-learning}

Every design so far is fixed by its designer. A \textbf{learning agent} improves
itself through four components:
\begin{itemize}
  \item the \textbf{policy} chooses an action, given the context (the percepts
    and the internal state);
  \item the \textbf{critic} evaluates the current behavior against an external
    performance measure and sends feedback;
  \item the \textbf{learning element} uses the feedback to change the other
    components;
  \item the \textbf{problem generator} steers the agent away from greedy
    behavior by trying out new actions, which balances the critic.
\end{itemize}
Learning is orthogonal to the previous types: any of them can be a learning
agent or not, and learning can modify any of its components.

\subsection{Knowledge representation}\label{sec:agents-representation}

% ── slides L2 p.29–31 (diagrams) ──
Every agent above represents the world somehow. The lecture places a
representation on two axes.

\textbf{Expressiveness.} An \textbf{atomic} representation treats each state as
indivisible, with no internal structure. A \textbf{factored} representation
describes a state as a vector of attribute values. A \textbf{structured}
representation describes a state as a graph: objects with their properties and
the relationships among them. Expressiveness grows from atomic to structured.

\textbf{Distribution.} In a \textbf{localist} representation, concepts and
memory locations correspond one to one. In a \textbf{distributed}
representation, one memory cell takes part in several concepts, and one
concept is spread over several cells.
